\clearpage
\SourceLeaf{172}\centerline{[ZUR THEORIE DER PARALLELLINIEN.]}

\bigskip
\hrule
\bigskip

\centerline{[1.]}

\smallskip

\centerline{\textsc{Gauss} an \textsc{Olbers}. \ Braunschweig, 30. Juli 1806.}

\smallskip

.\ .\ .\ .\ .\ Hr. \textsc{von Zach} schreibt mir noch, dass Sie Sich zur Recension von \textsc{Legendres} Werk über die Kometenbahnen erboten hätten. Mit Vergnügen werde ich Ihnen also das mir von Hrn. \textsc{von Zach} zugeschickte Exemplar nach Bremen senden, doch erlauben Sie wohl, dass ich es erst noch einige Wochen behalte. Bei vorläufigem Durchblättern scheint es mir sehr viel Schönes zu enthalten. Vieles von dem, was ich in meiner Methode besonders in ihrer ersten Gestalt Eigenthümliches hatte, finde ich auch in diesem Buche wieder. Es scheint mein Schicksal zu sein, fast in allen meinen theoretischen Arbeiten mit \textsc{Legendre} zu concurriren. So in der höhern Arithmetik, in den Untersuchungen über transcendente Functionen, die mit der Rectification der Ellipse zusammenhangen, bei den ersten Gründen der Geometrie und nun wieder hier. So ist z.\,B. auch das von mir seit 1794 gebrauchte Princip, dass man, um mehrere Grössen, die man nicht alle genau darstellen kann, am besten darzustellen, die Summe der Quadrate zu einem Minimum machen müsse, auch in \textsc{Legendres} Werke gebraucht und recht wacker ausgeführt.\ .\ .\ .\ .\ .\ .

\bigskip
\hrule
\bigskip

\centerline{[2.]}

\smallskip

\centerline{[\emph{Aus} \textsc{Schumachers} \emph{Tagebuch: »Gaussiana«. Göttingen, November 1808.}]}

\smallskip

\{\textsc{Gauss} hat die Theorie der Parallellinien darauf zurückgebracht, dass wenn die angenommene Theorie nicht wahr wäre, es eine constante a priori der Länge nach gegebene Linie geben müsste, welches absurd ist. Doch hält er selbst diese Arbeit noch nicht für hinreichend.\}

\bigskip
\hrule
